\section{From passive listening to biofeedback}

In the last section it was established how heart rate data is extracted using rPPG techniques. The next step is to analyse how this data can be translated into a piece of music. It is first necessary to review the existing research on passive music listening and its limitation in stress reduction.

According to a meta-analysis by Pelletier (2004), music consistently assists in lowering stress levels. The magnitude of this response is affected by several factors. Patients with prior musical experience, such as playing an instrument, tend to lower their stress levels a lot more, as they often possess a greater predisposition and comfort level with music \cite{pelletier2004effect}.
Furthermore, Pelletier notes that individual therapy generally results in better relaxation than group session. Although, group sessions are used much more often because they are less expensive. A huge obstacle is that the devices that are used for data collection often cause discomfort. The study highlights that physiological measures or self-reports can cause a disruption of the relaxed state. Behavioral observation seems to collect more accurate data of the the relaxation process  \cite{pelletier2004effect}.
Moreover, the study came to the conclusion that music based on scientific research yiels greater effects for stress reduction than music that was chosen on personal preference  \cite{pelletier2004effect}.
This is likely because preferred music can be too distracting or stimulating. On the other hand, research-based music combined with techniques like verbal suggestion or progressive relaxation achieves the highest impact on decreasing arousal. \cite{pelletier2004effect}.
In another study, Linnemann et al. examined how music affects stress in daily life rather than in a laboratory setting. They concluded that while music can effectively reduce stress levels, the magnitude of this effect depends heavily on the listener's intent. Specifically, the largest reduction is achieved when subjects listen with the explicit goal of decreasing stress. This suggests that the underlying motivation is more important for the actual relaxation process than the music itself. Conversely, the study confirms that using music for the purpose of distraction can actually lead to an increase in stress levels \cite{linnemann2015music}.


The study of Bernardi et al. (2006) confirms that the results do not depend on the music preferation of the user rather on the BPM and rhythm structure of the music. Additionally, the study found no evidence of habituation, meaning the users do not simply 'get used' to the music over time.
Their findings suggest that because the body snychronizes with the music's tempo and rhythm measurable stress-reducing results can occur during passive listening, regardless of the users' conscious intention \cite{bernardi2006cardiovascular}.

% hier drunter noch einmal alles checken.

